\documentclass[11pt, a4paper, oneside]{book}
\usepackage{titling}
\title{Heading 1}
\date{2015-2016}
\author{Lorem Ipsum Author}
\usepackage{helvet} 				%Helvetica Font
\usepackage{helvet} 				% Helvetica Font
\usepackage{baskervald}			% Baskervald font
\usepackage{lipsum}				% Om Lorem Ipsum te maken
\usepackage{tabularx}			% Make tables
\usepackage{multicol}				% Create multiple columns
\usepackage{enumitem}			% Make Lists
\usepackage{listings}				% Display code, with syntax coloring and indenting
\usepackage[dutch]{babel}		% Nederlandse tekenafstand - woordafbraak
\usepackage{mdframed}			% Make text boxes
\usepackage[usenames,dvipsnames]{xcolor}
\usepackage{color}				% Kleuren definieren met variabelen
\usepackage{hyperref}			% Printing hyperlinks
\usepackage{titlesec}				% Enhancement/replacement of the standard LaTeX title, headers and content macros
\usepackage{url}				% Format url's in a better way, is not a replacement for the hyperref package
\usepackage{amsmath}			% Enhancements for printing math formulas
\usepackage{eurosym}			% Use the euro symbol in LaTeX

% Om afbeeldingen te importeren

	\usepackage{graphicx}
	\graphicspath{ {images/} }

% Kleuren

	% Later systeem bedenken om ook CMYK colors te doen.
	\definecolor{grijs}{RGB}{150,150,150}
	\definecolor{oranjerood}{RGB}{240,76,37}
	\definecolor{oranjerood50pc}{RGB}{248,163,127}
	\definecolor{blauwgroen}{RGB}{0,156,171}
	\definecolor{blauwgroen50pc}{RGB}{111,195,206}

% HUISSTIJL ACHTERGROND VOORPAGINA

	\usepackage[pages=some]{background}

% With the eso-pic package, pages can have backgrounds, while still having text on top, this template used this package to add the thomas more background to the title page.
\usepackage{eso-pic}
\newcommand\BackgroundPic{%
	\put(0,0){%
		\parbox[b][\paperheight]{\paperwidth}{%
			\vfill
			\centering
			\includegraphics[width=\paperwidth,height=\paperheight,%
			keepaspectratio]{images/ThomasMoreBackground.png}%
			\vfill
			}
		}
	}

% ============

% Te veel subsection ruimte oplossen

	\usepackage{titlesec}
	\titlespacing{\subsection}{0pt}{*1}{*1}

% ===========

% Font veranderen

	\renewcommand*{\familydefault}{\sfdefault}

% ===========

\usepackage[utf8]{inputenc}
\usepackage{lmodern} % load a font with all the characters

\pagestyle{plain} % Alleen een page number wordt getoond

% Edit the margins
\geometry{a4paper,tmargin=.5in,bmargin=.5in,lmargin=.5in,rmargin=.7in,headheight=0in,headsep=0in,footskip=.3in}
 
\setlength{\parindent}{.1in}
\setlength{\parskip}{.1in}

\setlist[itemize]{itemsep=0mm}

\renewcommand*\contentsname{Inhoud}

%%%%%%%%%%%%

% Hoofdstuk tekst verwijderen en voorkomen
% dat Latex een nieuwe pagina maakt voor een nieuw hoofdstuk.

\usepackage{etoolbox} % For classes and package authors, to edit macro below

\makeatletter
\patchcmd{\chapter}{\if@openright\cleardoublepage\else\clearpage\fi}{}{}{}
\makeatother

%\titleformat{\chapter}{\normalfont\huge}{\thechapter.}{20pt}{\huge\it}
%
%\titleformat{\chapter}[display]
%    {\normalfont\large\bfseries}{}{2pt}{\Huge}
%\titlespacing*{\chapter}{0pt}{-40pt}{10pt}
%
%\setlength{\parindent}{0ex}

%%%--- MAIN DOCUMENT ---%%%

% Begin The document
\begin{document}
% VOORPAGINA
	\begin{titlepage}
			% Code is nog niet klaar, later variabelen toevoegen om alles makkelijker te kunnen aanpassen en wat betere code gebruiken.

				 \AddToShipoutPicture*{\BackgroundPic}

				\pagenumbering{gobble}
				\vspace*{130mm}
				\begin{flushright}
				\bf
				\color{oranjerood}
				{
					\scshape	
					\huge
					\fontfamily{phv}
					\selectfont
					\thetitle
				}
				
				\vspace{0mm}
				\LARGE
				Heading 2
				\end{flushright}
				\vspace{\fill}
				{\color{blauwgroen} Heading 3}

				\newpage

	\end{titlepage}

% INHOUDSTAFEL

\tableofcontents

\newpage

\pagenumbering{arabic}

% BEGIN


\chapter{Inleiding}

Dit is een LaTeX template voor een Thomas More hogeschool paper.
 \includegraphics[scale=0.5]{images/ThomasMoreBackground.png}

\newpage

\chapter{Winst}

	\begin{mdframed}[backgroundcolor=grijs!40,shadow=false,roundcorner=8pt]
		$$test = test - test$$
	\end{mdframed}

\chapter{Testhoofdstuk 1}

 Cras condimentum vestibulum tortor id venenatis. Suspendisse ac velit eget lacus volutpat egestas id eget nibh. Maecenas eget venenatis eros. Duis luctus pharetra est, ac consequat quam hendrerit vel. Suspendisse in aliquam lectus. Nunc malesuada mi a mauris posuere, eu aliquam augue dignissim. Vestibulum faucibus ipsum id orci rhoncus facilisis. Nulla facilisi. Aliquam vitae faucibus ipsum, at ultricies lacus. Aenean posuere mauris dolor, ac posuere tortor rhoncus id.

	\begin{mdframed}[backgroundcolor=grijs!40,shadow=false,roundcorner=8pt]
		$$Testhoofdstuk_A = Testhoofdstuk_B - Testhoofdstuk_C$$
	\end{mdframed}

\end{document}
